\documentclass[10pt,a4paper,oneside]{article}

\usepackage[T1]{fontenc}
\usepackage[utf8]{inputenc}
\usepackage{lmodern}
\usepackage[a4paper,margin=2.2cm,headheight=20pt]{geometry}
\usepackage{microtype}
\usepackage{booktabs}
\usepackage{longtable}
\usepackage{tabularx}
\usepackage{array}
\usepackage{enumitem}
\usepackage{xcolor}
\usepackage{listings}
\usepackage{fancyhdr}
\usepackage{titlesec}
\usepackage{amsmath}
\usepackage{hyperref}

\definecolor{navy}{HTML}{17365D}
\definecolor{blue}{HTML}{1F5A94}
\definecolor{ink}{HTML}{1F2933}
\definecolor{muted}{HTML}{52606D}
\definecolor{line}{HTML}{CBD5E1}
\definecolor{pale}{HTML}{F2F6FA}
\definecolor{codebg}{HTML}{F7F8FA}
\definecolor{green}{HTML}{176B4D}
\definecolor{amber}{HTML}{8A5A00}

\hypersetup{
  colorlinks=true,
  linkcolor=blue,
  urlcolor=blue,
  citecolor=blue,
  pdftitle={Reversede Technical Reference and Operating Guide},
  pdfauthor={Guilherme Moura Fernandes}
}

\setlist{itemsep=2pt,topsep=4pt}
\setlength{\parindent}{0pt}
\setlength{\parskip}{5pt}
\titleformat{\section}{\Large\bfseries\color{navy}}{\thesection}{0.7em}{}
\titleformat{\subsection}{\large\bfseries\color{blue}}{\thesubsection}{0.7em}{}
\titleformat{\subsubsection}{\normalsize\bfseries\color{ink}}{\thesubsubsection}{0.7em}{}

\pagestyle{fancy}
\fancyhf{}
\lhead{\small\textcolor{muted}{Reversede \textbar{} Technical Reference}}
\rhead{\small\textcolor{muted}{v0.3.0.dev0}}
\cfoot{\small\textcolor{muted}{\thepage}}
\renewcommand{\headrulewidth}{0.4pt}
\renewcommand{\headrule}{\hbox to\headwidth{\color{line}\leaders\hrule height \headrulewidth\hfill}}

\lstdefinestyle{reversede}{
  basicstyle=\ttfamily\scriptsize,
  backgroundcolor=\color{codebg},
  frame=single,
  rulecolor=\color{line},
  framerule=0.4pt,
  xleftmargin=5pt,
  xrightmargin=5pt,
  aboveskip=7pt,
  belowskip=7pt,
  breaklines=true,
  breakatwhitespace=true,
  columns=fullflexible,
  keepspaces=true,
  showstringspaces=false,
  numbers=left,
  numberstyle=\tiny\color{muted},
  numbersep=7pt
}
\lstset{style=reversede}

\newcommand{\code}[1]{\texttt{\small\detokenize{#1}}}
\newcommand{\note}[1]{\textcolor{blue}{\textbf{Note.} #1}}
\newcommand{\warning}[1]{\textcolor{amber}{\textbf{Boundary.} #1}}
\newcommand{\ok}[1]{\textcolor{green}{\textbf{Verified.} #1}}
\newcommand{\labelrow}[2]{\textbf{#1} & #2\\}

\begin{document}

\begin{titlepage}
\thispagestyle{empty}
\vspace*{1.4cm}
{\color{navy}\rule{\textwidth}{2pt}}\\[1.0cm]
{\Huge\bfseries\color{navy} Reversede}\par
\vspace{0.25cm}
{\LARGE\bfseries Technical Reference and Operating Guide}\par
\vspace{0.35cm}
{\large A documentation reference for auditable DataFrame diagnostics and scikit-learn blueprints}\par
\vspace{1.0cm}
\begin{tabularx}{\textwidth}{>{\bfseries\color{muted}}p{4.2cm}X}
Package distribution & \code{reversede-ai}\\
Public import & \code{reversede}\\
Current package version & \code{0.3.0.dev0}\\
Python support & CPython 3.10-3.13\\
Source policy & Private GitHub repository, access controlled by the owner\\
Proposed author attribution & Guilherme Moura Fernandes\\
Document status & Technical reference for the private development line\\
Date & 12 October 2026
\end{tabularx}
\vfill
\fcolorbox{navy}{pale}{\begin{minipage}{0.93\textwidth}
\textbf{Scope.} This manual documents the public API, algorithmic logic, generated artifacts,
portability model, and reproducibility boundaries of the current library. It is written as a
reference document rather than as a research claim. Historical study results were produced with
an earlier 0.2.0 line and are not automatically attributed to this release.
\end{minipage}}
\vfill
{\color{navy}\rule{\textwidth}{2pt}}\\[0.35cm]
\small\textcolor{muted}{Prepared from the private repository source and its versioned documentation.}
\end{titlepage}

\tableofcontents
\newpage

\section{Executive summary}

Reversede is a Python library that turns a pandas \code{DataFrame} into an auditable analysis
surface: it profiles columns, resolves or questions a target, detects structural and data-quality
stress, ranks candidate features, writes model-conditioned prescriptions, and emits an unfitted
scikit-learn preprocessing blueprint. The central design choice is to expose evidence and uncertainty
instead of hiding decisions inside a single opaque score.

The library does not train an estimator, does not select a winning model, and does not establish
causality, absence of leakage, or superiority over a baseline. A generated blueprint is source code
that still needs to be reviewed, fitted on training folds, evaluated on held-out data, and adapted to
the user's data contract.

\subsection{How to read this reference}
\begin{enumerate}
  \item Start with Section~\ref{sec:contract} for the package contract and installation paths.
  \item Use Sections~\ref{sec:algorithm}-\ref{sec:api} when implementing or reviewing code.
  \item Use Sections~\ref{sec:repro}-\ref{sec:portability} when preparing a reproducible run.
  \item Use the troubleshooting and glossary sections as an operational checklist.
\end{enumerate}

\section{Contract at a glance}\label{sec:contract}

\begin{tabularx}{\textwidth}{>{\bfseries}p{4.1cm}X}
\toprule
Item & Contract\\
\midrule
Installable distribution & \code{reversede-ai}\\
Import namespace & \code{from reversede import Reversede, ReversedeConfig}\\
Runtime & Python 3.10 or newer; the maintained matrix covers 3.10, 3.11, 3.12, and 3.13\\
Required dependencies & \code{numpy>=1.26}, \code{pandas>=2.1}\\
Optional ML dependency & \code{scikit-learn>=1.4} via \code{[sklearn]} or \code{[ml]}\\
Optional notebook dependency & \code{jupyterlab>=4}, \code{ipykernel>=6}, scikit-learn\\
Primary input & A non-empty pandas \code{DataFrame}\\
Primary output & Evidence-bearing object plus JSON-ready summary and source blueprint\\
Mutation policy & The caller's DataFrame is not mutated; analysis works on a deep copy\\
Distribution policy & The source repository is private; collaborators require explicit access\\
\bottomrule
\end{tabularx}

\subsection{The smallest useful call}
\begin{lstlisting}[language=Python]
import pandas as pd
from reversede import Reversede

frame = pd.read_csv("data.csv")
analysis = Reversede(frame, target="target", task="regression")
print(analysis.summary())
source = analysis.generate_blueprint("artifacts/preprocess.py")
\end{lstlisting}

The explicit target and task are recommended when the data contract is known. Inference remains
available for exploratory work, but it reports confidence, candidates, and abstention signals.

\section{Installation and access}

The repository is private by design. A person may install or clone the project only after the owner
has granted GitHub access. The following commands assume an authenticated Git client or a configured
credential helper; no token is embedded in source, notebooks, or this manual.

\subsection{Virtual environment and direct Git installation}
\begin{lstlisting}[language=bash]
python -m venv .venv
# Windows PowerShell
.\.venv\Scripts\Activate.ps1
# macOS or Linux
source .venv/bin/activate

python -m pip install --upgrade pip
python -m pip install "reversede-ai @ git+https://github.com/alemofer-hash/reversede-ai.git@v0.3.0-dev0"
\end{lstlisting}

The direct URL resolves a private Git ref. Plain \code{pip install reversede-ai} is not the selected
distribution path for this private development line because it searches a package index rather than
the controlled repository.

\subsection{Extras}
\begin{longtable}{>{\ttfamily}p{3.1cm}p{4.2cm}p{7.1cm}}
\toprule
Extra & Main packages & Intended use\\
\midrule
\endhead
\code{sklearn}, \code{ml} & scikit-learn & Blueprint execution and preprocessing pipelines\\
\code{profiling}, \code{report} & ydata-profiling, setuptools<81 & Optional HTML profile adapter\\
\code{analysis} & matplotlib, scipy, seaborn, scikit-learn & Notebook analysis and visual diagnostics\\
\code{notebook} & JupyterLab, ipykernel, scikit-learn & Jupyter, Colab, and notebook kernels\\
\code{test} & pytest & Local test execution\\
\code{dev} & build, pytest, pytest-cov, ruff, twine & Maintainer quality and packaging checks\\
\bottomrule
\end{longtable}

\subsection{IDE and notebook setup}
The same Python environment can be selected by VS Code, PyCharm, JupyterLab, or Anaconda. The
repository includes an \code{environment.yml}, pinned requirement families, VS Code settings, a
dev-container definition, and two synthetic notebooks. The reliable sequence is:

\begin{enumerate}
  \item create or select one environment;
  \item install the private Git ref and the required extra;
  \item register that environment as a Jupyter kernel;
  \item run the import smoke check below before opening a long notebook.
\end{enumerate}

\begin{lstlisting}[language=Python]
import sys
import reversede

print(sys.executable)
print(reversede.__version__)
\end{lstlisting}

In Google Colab, the same authenticated Git URL can be used in a cell after access has been granted.
Credentials must be supplied through the notebook runtime or a secret manager and must never be
written into a committed cell.

\section{Architecture and data flow}\label{sec:architecture}

The public class is a thin orchestration boundary. Each internal module has one reviewable
responsibility, and the intermediate structures are retained on the analysis object.

\begin{center}
\fcolorbox{navy}{pale}{\begin{minipage}{0.94\textwidth}
\centering
\textbf{DataFrame + target/task/config}\par\smallskip
$\Downarrow$\par\smallskip
\fbox{validate input} $\rightarrow$ \fbox{profile and practical types} $\rightarrow$
\fbox{resolve target and task}\par\smallskip
$\Downarrow$\par\smallskip
\fbox{association and eligibility} $\rightarrow$ \fbox{stress diagnostics} $\rightarrow$
\fbox{model-conditioned prescriptions}\par\smallskip
$\Downarrow$\par\smallskip
\fbox{ranked candidates + evidence + alerts + unfitted sklearn blueprint}
\end{minipage}}
\end{center}

\begin{longtable}{>{\ttfamily}p{3.1cm}p{5.2cm}p{6.1cm}}
\toprule
Module & Responsibility & Review output\\
\midrule
\endhead
\code{validation.py} & Validate arguments, shape, duplicate names, and target & Fail-fast errors with actionable messages\\
\code{types.py} & Infer practical type and signals for IDs, text, dates, and numeric strings & Type evidence used by downstream decisions\\
\code{target.py} & Resolve explicit or inferred target, confidence, margin, and abstention & Target source, confidence, candidates, alerts\\
\code{association.py} & Compute Pearson/Spearman, eta squared, and Cramer V where valid & Comparable feature-target evidence\\
\code{profile.py} & Build column profile and eligibility explanations & Profile and candidate rows\\
\code{stress.py} & Detect missingness, infinities, outliers, scale, cardinality, dates, and collinearity & Stress records with severity and counts\\
\code{prescriptions.py} & Convert evidence into model-conditioned recommendations & Action, rationale, affected models, confidence\\
\code{blueprint.py} & Generate syntactically checked sklearn source without fitted global values & Reproducible preprocessing source\\
\code{api.py} & Orchestrate the pipeline and persist reports & Stable public object and report files\\
\bottomrule
\end{longtable}

\section{Algorithm and decision logic}\label{sec:algorithm}

\subsection{End-to-end pseudocode}
\begin{lstlisting}
function REVERSEDE(data, target, task, config):
    validate_dataframe_and_arguments(data, target, task, config)
    working_copy <- deep_copy(data)
    profile <- profile_all_columns(working_copy, config)

    if target is explicit:
        resolved_target <- validate_and_accept(target)
        target_source <- "explicit"
        target_confidence <- 1.0
    else:
        candidates <- score_target_candidates(profile, working_copy)
        resolved_target, confidence <- infer_with_margin(candidates, config)
        target_source <- "inferred" or "abstained"

    if task is explicit:
        resolved_task <- task
        task_source <- "explicit"
    else:
        resolved_task <- resolve_task_from_target(profile, working_copy)
        task_source <- "heuristic"

    for feature in profile excluding resolved_target:
        explain_eligibility(feature, config)
        if target is reliable:
            compute_compatible_association(feature, resolved_target, resolved_task)
            compute_candidate_score(feature)
        else:
            compute_structural_priority_only(feature)

    candidates <- rank_and_limit_features(config.max_features)
    stress <- diagnose_stress(working_copy, profile, candidates, config)
    prescriptions <- build_model_conditioned_prescriptions(profile, stress, resolved_task)
    blueprint <- generate_unfitted_sklearn_source(profile, candidates, stress, config)
    return evidence_object(profile, target, candidates, stress,
                           prescriptions, blueprint, alerts)
\end{lstlisting}

\subsection{Target resolution}
An explicit target is accepted only if it names an existing, non-degenerate column. Without one,
the target resolver scores plausible columns and applies both a confidence threshold and a margin
over the runner-up. If the margin is insufficient, the library abstains. This is deliberate: an
incorrect target can make every downstream association look meaningful.

\begin{center}
\begin{tabular}{p{4.0cm}p{10.3cm}}
\toprule
State & Meaning\\
\midrule
Explicit & The caller owns the semantic decision; confidence is reported as 1.0 for provenance.\\
Inferred & A heuristic candidate passed configured threshold and margin checks.\\
Abstained & Evidence was insufficient or ambiguous; structural diagnostics remain available, but predictive importance is not claimed.\\
\bottomrule
\end{tabular}
\end{center}

\subsection{Feature eligibility and ranking}
The candidate table distinguishes eligibility from usefulness. Identifier-like columns, constant
columns, unsupported objects, and the resolved target receive reasons rather than silently entering
the shortlist. When a target is reliable, numeric and categorical associations are calculated with
compatible measures. When it is not, the library falls back to structural priority and labels it as
such.

Scores are heuristic ranking aids on a normalized scale. They are not probabilities, causal effects,
feature importance from a fitted estimator, or guarantees of generalization.

\subsection{Stress and prescriptions}
Stress diagnostics quantify the conditions that can break a downstream model: missing and infinite
values, near constants, outliers, scale imbalance, high cardinality, rare categories, date-like
columns, and correlated numeric pairs. A prescription keeps the evidence, proposed action, affected
model families, and confidence together. The user can therefore reject or revise a prescription
without losing the observation that caused it.

\section{Public API}\label{sec:api}

\subsection{Constructor}
\begin{lstlisting}[language=Python]
Reversede(
    data: pandas.DataFrame,
    target: str | None = None,
    task: str = "auto",
    config: ReversedeConfig | None = None,
)
\end{lstlisting}

The constructor performs validation, profiling, target/task resolution, candidate analysis, stress
diagnostics, and prescription generation. It does not fit a model. The input frame is copied before
analysis; the original object remains under caller control.

\subsection{Methods}
\begin{longtable}{>{\ttfamily}p{5.4cm}p{8.8cm}}
\toprule
Method & Behavior\\
\midrule
\endhead
\code{generate_blueprint(path=None, print_code=False) -> str} & Return a checked sklearn preprocessing source string. If \code{path} is given, persist the source; if \code{print_code} is true, also print it. No estimator is selected or trained.\\
\code{summary() -> dict[str, Any]} & Return JSON-ready, evidence-bearing state for reports, review, or a notebook cell.\\
\code{save_report(directory) -> dict[str, str]} & Write CSV/JSON summaries and the generated blueprint under the requested directory and return their paths.\\
\bottomrule
\end{longtable}

\subsection{Stable fields}
\begin{tabularx}{\textwidth}{>{\ttfamily}p{4.2cm}X}
\toprule
Field & Interpretation\\
\midrule
\code{target_}, \code{target_source_}, \code{target_confidence_}, \code{target_candidates_} & Resolved target and the evidence path that produced it.\\
\code{task_}, \code{task_source_} & Regression/classification decision and whether it was explicit or heuristic.\\
\code{profile_} & Per-column types, counts, missingness, cardinality, and signals.\\
\code{candidates_} & Eligibility, associations, structural priority, and shortlist scores.\\
\code{stress_} & Stress diagnostics with values, thresholds, severity, and affected features.\\
\code{prescriptions_} & Conditional recommendations with rationale and model families.\\
\code{alerts_} & Warnings for ambiguity, leakage risk, unsupported shape, or review conditions.\\
\code{selected_features_} & The final shortlist after configured limits and exclusions.\\
\bottomrule
\end{tabularx}

\section{Parameterization and configuration}

\code{ReversedeConfig} is a frozen dataclass. Its values are part of the analysis contract and are
propagated into the generated blueprint where applicable.

\begin{longtable}{>{\ttfamily}p{4.3cm}p{2.1cm}p{7.8cm}}
\toprule
Parameter & Default & Effect and review question\\
\midrule
\endhead
\code{max_features} & 20 & Maximum shortlist size. Is the downstream model allowed to see more or fewer columns?\\
\code{max_categorical_levels} & 50 & Upper bound used by categorical blueprint handling. Does the chosen bound fit memory and business semantics?\\
\code{high_cardinality_levels} & 25 & Warning boundary for categorical cardinality. Is the column an ID, a real category, or a text field?\\
\code{near_constant_threshold} & 0.995 & Dominant-value fraction that triggers a near-constant warning.\\
\code{id_unique_fraction} & 0.98 & Uniqueness fraction that signals a likely identifier.\\
\code{target_confidence_threshold} & 0.88 & Minimum confidence for accepting an inferred target.\\
\code{target_margin} & 0.08 & Required gap between the best and second target candidate.\\
\code{rare_category_fraction} & 0.01 & Frequency boundary for rare-category warnings and prescriptions.\\
\code{outlier_fraction_threshold} & 0.05 & Fraction of observations above the outlier diagnostic boundary.\\
\code{include_columns} & \code{()} & Explicitly include columns otherwise excluded; the reason remains visible.\\
\code{exclude_columns} & \code{()} & Explicitly exclude columns from the shortlist and blueprint.\\
\bottomrule
\end{longtable}

\begin{lstlisting}[language=Python]
from reversede import Reversede, ReversedeConfig

config = ReversedeConfig(
    max_features=30,
    include_columns=("business_signal",),
    exclude_columns=("record_id",),
)
analysis = Reversede(frame, target="target", task="classification", config=config)
config.validate()
\end{lstlisting}

Invalid ranges are rejected by \code{validate()}. Explicit inclusions remain auditable rather than
being converted into an unmarked override.

\section{Generated blueprint and pseudocode for use}

\subsection{What the blueprint contains}
The generated source is designed to be a readable preprocessing starting point. Depending on column
signals, it may include numeric imputation and scaling, conversion of numeric strings, categorical
imputation and one-hot encoding, and datetime handling. Categorical limits and frequency thresholds
come from the configuration. The output is intentionally unfit: it contains transformer definitions,
not learned medians, means, category vocabularies, or a fitted model.

\begin{lstlisting}[language=Python]
blueprint = analysis.generate_blueprint(
    path="artifacts/reversede_preprocessing.py",
    print_code=False,
)
\end{lstlisting}

\subsection{Fit only inside the training boundary}
\begin{lstlisting}[language=Python]
from sklearn.model_selection import train_test_split
from sklearn.pipeline import Pipeline
from sklearn.linear_model import LogisticRegression

train, test = train_test_split(frame, test_size=0.2, random_state=42)
# Import or adapt the generated preprocessing source here.
preprocessor = build_preprocessor_from_blueprint(train)
model = Pipeline([
    ("preprocess", preprocessor),
    ("estimator", LogisticRegression(max_iter=1000)),
])
model.fit(train.drop(columns="target"), train["target"])
score = model.score(test.drop(columns="target"), test["target"])
\end{lstlisting}

This pattern prevents test-set statistics from leaking into preprocessing. The blueprint does not
prove that the eventual training code follows this rule; the experiment owner remains responsible.

\subsection{Optional profiling adapter}
\begin{lstlisting}[language=Python]
from reversede.integrations.profiling import profile_with_ydata

report = profile_with_ydata(frame, "artifacts/profile.html", minimal=True)
\end{lstlisting}

The adapter is optional and raises an actionable \code{ImportError} when its extra is absent.

\section{Report and output schemas}

\subsection{Summary shape}
The exact internal dictionaries may grow, but the following conceptual shape is stable and intended
for JSON serialization:

\begin{lstlisting}
{
  "version": "0.3.0.dev0",
  "target": {"name": "target", "source": "explicit", "confidence": 1.0},
  "task": {"name": "regression", "source": "explicit"},
  "profile": [...],
  "candidates": [...],
  "stress": [...],
  "prescriptions": [...],
  "alerts": [...],
  "selected_features": [...]
}
\end{lstlisting}

\subsection{Persisted report}
\code{save_report(directory)} returns a dictionary of paths for the serialized summary, tabular
candidate/stress artifacts, and the blueprint source. Treat these files as analysis outputs tied to
the input data, configuration, package version, and Python environment. A report is not a substitute
for a dataset manifest or a model evaluation record.

\section{Legacy adapter and migration}

The historical namespace \code{reverseda_ai} remains as a narrow adapter for 0.1.x callers. Its
\code{ReversedaAutonoma} class emits a \code{DeprecationWarning}. New code should use the English
namespace and methods.

\begin{longtable}{>{\ttfamily}p{5.2cm}>{\ttfamily}p{5.2cm}p{4.0cm}}
\toprule
Legacy & Current & Migration note\\
\midrule
\endhead
\code{reverseda_ai} & \code{reversede} & Import the public namespace.\\
\code{ReversedaAutonoma} & \code{Reversede} & Construct the new class directly.\\
\code{requisitos_} & \code{candidates_} & Use the richer, serializable candidate surface.\\
\code{estresse_} & \code{stress_} & Includes additional diagnostics and evidence.\\
\code{gerar_blueprint_extrator()} & \code{generate_blueprint()} & English name and explicit persistence controls.\\
\code{resumo()} & \code{summary()} & JSON-ready report surface.\\
\bottomrule
\end{longtable}

\section{Reproducibility, provenance, and security}\label{sec:repro}

\subsection{Minimum reproducibility record}
For a repeatable run, store the following beside the report:
\begin{itemize}
  \item package version and Git ref;
  \item Python executable and operating system;
  \item dependency lock or requirement file;
  \item configuration values and explicit target/task;
  \item dataset provenance, schema, row/column counts, and a content hash where permitted;
  \item command or notebook cell used to create the report;
  \item blueprint source and the downstream estimator/evaluation protocol.
\end{itemize}

\subsection{Private repository policy}
The GitHub repository \code{alemofer-hash/reversede-ai} is private. Access is granted by the owner
to selected collaborators. A private repository is an access control boundary, not proof of authorship,
license clearance, or data provenance. Before sharing a clone or an artifact, review:

\begin{enumerate}
  \item authorship and contribution records;
  \item licenses and notices for third-party dependencies and copied material;
  \item notebooks, PDFs, datasets, personal paths, and generated outputs;
  \item secrets, tokens, certificates, and local configuration;
  \item the intended audience and the minimum GitHub permission needed.
\end{enumerate}

\warning{Public release of code, notebooks, article, data, or descriptions of use is user-controlled
and outside this document's automatic workflow.}

\subsection{Scientific boundaries}
The library reports warnings and recommendations. It does not turn a finite diagnostic into a proof,
does not establish causal interpretation, and does not claim that one model family is superior.
The historical study used Reversede 0.2.0; this 0.3.0.dev0 reference documents a newer library line,
and the study was not rerun merely by publishing this manual.

\section{Portability matrix}\label{sec:portability}

\begin{longtable}{p{3.1cm}p{4.1cm}p{7.2cm}}
\toprule
Environment & Supported route & Operational note\\
\midrule
\endhead
VS Code & Select the project virtualenv or dev container & The repository includes editor settings and a container definition.\\
PyCharm & Configure an existing virtualenv or Conda interpreter & Run the import smoke check from the configured interpreter.\\
JupyterLab & Install \code{reversede-ai[notebook]} & Register and select the same environment as the kernel.\\
Anaconda & Use \code{environment.yml} or a Conda env plus pip Git install & Keep one authoritative interpreter for notebook and terminal.\\
Google Colab & Authenticated Git install in a cell & Use runtime secrets; never hard-code a private token.\\
CI & Matrix on Windows, Linux, macOS and Python 3.10-3.13 & Core, quality/build, and notebook smoke jobs are separated.\\
\bottomrule
\end{longtable}

\section{Verification evidence}

The private development line has a verification workflow proportionate to the library's scope:
\begin{itemize}
  \item unit test suite: 23 tests passed in the prepared local environment;
  \item local coverage measurement: 92 percent in the recorded run;
  \item Ruff quality checks passed;
  \item wheel and source distribution built and installed in an isolated target;
  \item synthetic notebooks executed after correcting a kernel-name mismatch and a notebook column reference;
  \item CI tag run passed across its configured matrix at the private GitHub repository.
\end{itemize}

These are software verification facts for the recorded environment. They are not evidence of
scientific superiority, universal data quality, or correctness on an arbitrary private dataset.

\section{Troubleshooting}

\begin{longtable}{p{4.1cm}p{10.3cm}}
\toprule
Symptom & First checks\\
\midrule
\endhead
Private install returns 404 & Confirm that the GitHub account has access and that the credential helper is active. Anonymous HTTP access is expected to fail.\\
Import works in terminal but not notebook & Compare \code{sys.executable} in both contexts and select the installed kernel.\\
Target is \code{None} & Inspect \code{target_candidates_}, confidence, margin, and alerts; pass an explicit target when the contract is known.\\
Too few candidates & Review ID/constant/unsupported-type reasons, exclusions, and \code{max_features}.\\
Blueprint output is not a model & This is intentional. Choose and validate an estimator separately, then fit preprocessing only on training folds.\\
Optional profiling import fails & Install the \code{profiling} or \code{report} extra in the active environment.\\
Results differ across machines & Freeze Python/dependencies, configuration, input hash, Git ref, and random seeds in the run record.\\
\bottomrule
\end{longtable}

\section{Glossary}
\begin{description}[style=nextline,leftmargin=3.2cm]
  \item[Candidate] A feature row carrying eligibility, evidence, and a ranking signal.
  \item[Confidence] A heuristic decision strength reported for target resolution; not a probability of truth.
  \item[Stress] A measurable data or structural condition that may affect downstream modeling.
  \item[Prescription] A conditional, evidence-linked action suggestion for a model family.
  \item[Blueprint] Unfitted preprocessing source code that a user reviews and fits inside a training boundary.
  \item[Abstention] An explicit refusal to make a weak or ambiguous target inference.
  \item[Provenance] The record of where data, code, dependencies, and claims came from.
\end{description}

\section{Citation and version record}

For internal citation, identify the distribution and Git ref rather than citing a mutable working
directory:

\begin{lstlisting}
Guilherme Moura Fernandes. Reversede: Technical Reference and Operating Guide.
reversede-ai 0.3.0.dev0, private GitHub repository, 2026.
\end{lstlisting}

The repository also contains \code{CITATION.cff}, \code{AUTHORS.md}, licensing and provenance notes,
compatibility documentation, and the private installation guide. Keep those records synchronized when
the public API or version changes.

\vfill
\begin{center}
\color{navy}\rule{0.72\textwidth}{0.7pt}\\[0.3cm]
\small\textcolor{muted}{End of technical reference. The source code and access policy remain private.}
\end{center}

\end{document}
